\section{The translation}
\label{sec:translation}

The translation $\tr{\cdot}$ (\lean{Translation.translate},
\file{Translation/Flow.lean}) compiles each \Lang{} thread into
continuation-passing Steps \cite{appel1992cps}. It maps a program to one
Flow definition with three kinds of Step types (\lean{StepType}):
\begin{itemize}
\item \texttt{point code}: a \emph{program point}. The Step type is the code
  left in the current procedure; the Step runs the head statement of that
  code and moves to the program point of what follows.
\item \lean{start}: the entry Step. It publishes one token for each
  statically known mutex and moves to the first program point of \kw{main}.
\item \lean{timed}: one generic Step type for every \kw{sleep} and every
  timed \kw{select}.
\end{itemize}
The Step \emph{input} is the thread's state: its thread ID, a counter for
fresh names, its local variables, and its stack of call frames and
exception handlers. A \kw{GoTo} passes that state on. So a thread is, at any
moment, exactly one active Step execution, and the thread's ``program
counter and stack'' are that Step's type and input.

\paragraph{One Step per statement.}
Each statement is its own Step. A statement that does not block has no
\kw{WaitFor}; its \kw{Execute} takes two commits, \rulename{STEP-START} and
\rulename{STEP-COMPLETE-EXECUTE}. A blocking statement's Step has a
\kw{WaitFor} that returns the wait, which adds two commits. Sequencing and
\kw{skip} are not Steps. These commit counts are the stamped actions of
Table~\ref{tab:timing}, and Figure~\ref{fig:lifecycle} shows the
correspondence for the two shapes.

\begin{figure}[t]
\centering
\begin{tikzpicture}[
  >={Stealth[length=4.5pt]},
  ph/.style={draw,rounded corners=2pt,font=\scriptsize\ttfamily,inner sep=3pt,minimum height=14pt,minimum width=1.5cm},
  dx/.style={draw,font=\scriptsize,inner sep=3pt,minimum height=14pt,minimum width=1.5cm,fill=black!6},
  lab/.style={font=\scriptsize},
  commit/.style={->,line width=1.1pt},
  other/.style={->},
  x=2.7cm,y=1cm]
% (a) emit
\node[lab,anchor=west] at (-0.6,3.75) {(a) \kw{emit}$(a)$: a Step without \kw{WaitFor}, two commits};
\node[lab,anchor=east] at (-0.34,3.0) {\Lang};
\node[ph] (l0) at (0,3.0) {ready};
\node[ph] (l1) at (1,3.0) {emitting};
\node[ph] (l2) at (3,3.0) {emitted};
\node[ph] (l3) at (5,3.0) {ready};
\draw[commit] (l0) -- node[lab,above] {begin} (l1);
\draw[other] (l1) -- node[lab,above] {$\mathit{out}(a)$} (l2);
\draw[commit] (l2) -- node[lab,above] {complete} (l3);
\node[lab,anchor=east] at (-0.34,1.9) {Dex};
\node[dx] (d0) at (0,1.9) {pending};
\node[dx] (d1) at (1,1.9) {Started};
\node[dx] (d2) at (2,1.9) {Running$[a]$};
\node[dx] (d3) at (3,1.9) {Running$[\,]$};
\node[dx] (d4) at (4,1.9) {Returned};
\node[dx] (d5) at (5,1.9) {next Step};
\draw[commit] (d0) -- node[lab,below] {start} (d1);
\draw[other] (d1) -- node[lab,below] {run} (d2);
\draw[other] (d2) -- node[lab,below] {$\mathit{fx}(a)$} (d3);
\draw[other] (d3) -- node[lab,below] {return} (d4);
\draw[commit] (d4) -- node[lab,below] {complete} (d5);
\foreach \a/\b in {l0/d0,l1/d1,l2/d3,l3/d5} \draw[densely dotted] (\a) -- (\b);
% (b) lock
\node[lab,anchor=west] at (-0.6,0.85) {(b) \kw{lock}$(m)$: a Step with a \kw{WaitFor}, four commits};
\node[lab,anchor=east] at (-0.34,0.1) {\Lang};
\node[ph] (b0) at (0,0.1) {ready};
\node[ph] (b1) at (1,0.1) {waitStarted};
\node[ph] (b2) at (2,0.1) {blocked};
\node[ph] (b3) at (3,0.1) {woken};
\node[ph] (b4) at (4,0.1) {resuming};
\node[ph] (b5) at (5,0.1) {ready};
\draw[commit] (b0) -- node[lab,above] {begin} (b1);
\draw[commit] (b1) -- node[lab,above] {register} (b2);
\draw[other,dashed] (b2) -- node[lab,above] {unlock} (b3);
\draw[commit] (b3) -- node[lab,above] {begin} (b4);
\draw[commit] (b4) -- node[lab,above] {complete} (b5);
\node[lab,anchor=east] at (-0.34,-1.0) {Dex};
\node[dx] (e0) at (0,-1.0) {$\mathit{WF}$, pending};
\node[dx] (e1) at (1,-1.0) {$\mathit{WF}$, task};
\node[dx] (e2) at (2,-1.0) {$\mathit{WT}$};
\node[dx] (e3) at (3,-1.0) {$\mathit{EX}$, pending};
\node[dx] (e4) at (4,-1.0) {$\mathit{EX}$, task};
\node[dx] (e5) at (5,-1.0) {next Step};
\draw[commit] (e0) -- node[lab,below] {start} (e1);
\draw[commit] (e1) -- node[lab,below] {complete} (e2);
\draw[other,dashed] (e2) -- node[lab,below] {token} (e3);
\draw[commit] (e3) -- node[lab,below] {start} (e4);
\draw[commit] (e4) -- node[lab,below] {complete} (e5);
\foreach \a/\b in {b0/e0,b1/e1,b2/e2,b3/e3,b4/e4,b5/e5} \draw[densely dotted] (\a) -- (\b);
\end{tikzpicture}
\caption{One \Lang{} statement and its Step. Upper lanes: the phases of the
thread (\lean{L.Phase}). Lower lanes: the Step's invocation and task. Thick
arrows are clock-stamped: a stamped \Lang{} action above, a Dex commit
(\rulename{STEP-START} or a Complete rule) below. (a) An \kw{emit} is
issued at \rulename{STEP-START}; $\ev{out}(a)$ is the Worker's effect; the
thread continues at \rulename{STEP-COMPLETE-EXECUTE}. The Worker's run and
return steps have no \Lang{} counterpart. (b) A blocked \kw{lock} is woken
by another thread's \kw{unlock} (dashed): the commit that publishes the
token moves the waiting Step from $\mathit{WT}$ to $\mathit{EX}$.}
\label{fig:lifecycle}
\end{figure}

\paragraph{The \kw{Execute} of a program point.}
\lean{pointStepDef} runs the thread-local semantics of the head statement
(\lean{L.completeStatement}, the same function \Lang{}'s own semantics
uses) and encodes the atomic action it asks for as Dex effects and a
decision (\lean{encodeIntent}). \lean{encodeIntent} is the mapping table of
the specification, one case per row; Table~\ref{tab:mapping} reproduces it.

\begin{table}[t]
\centering
\small
\begin{tabularx}{\textwidth}{@{}P{3.5cm}L@{}}
\toprule
\Lang{} construct & Dex mapping \\
\midrule
\kw{if}, \kw{while}, local assignment & A Step whose \kw{Execute} \kw{GoTo}s
  the next program point (the taken branch, the loop body, or the loop
  exit). \\
Shared read, shared write & One \kw{Execute} per access: read from the
  snapshot, or write one Attribute. \\
Call, \kw{return}, recursion & Push or pop a frame on the explicit stack in
  the Step input. No SubFlow. \\
\kw{throw}, \kw{try} & The handler stack is in the Step input; \kw{throw}
  is a \kw{GoTo} to the innermost handler's program point. No Worker
  error. \\
\kw{fork} & \kw{GoToMany}: the parent's continuation and the child's first
  program point. The child's thread ID is the parent's ID extended by the
  parent's counter. \\
Thread end & Publish one message to the Channel $\mathit{join}/t$, then
  \kw{DeadEnd}. \\
\kw{join}$(t)$ & \kw{WaitFor} returns
  $\kw{AnyOf}(\mathit{join}/t,\ \mathit{cancel}/\mathit{self})$. \\
\kw{exit}$(v)$, return of \kw{main} & \kw{ForceComplete} with the tagged
  value $\{\texttt{exit}: v\}$, which is never null. \\
\kw{fail}$(s)$, uncaught string exception & \kw{ForceFail}$(s)$. \\
Halted thread & \kw{DeadEnd}, publishing nothing. \\
\kw{newMutex}, \kw{unlock}$(m)$ & \kw{Execute} publishes one token to
  $m$'s token Channel. \\
\kw{lock}$(m)$ & \kw{WaitFor} returns $\kw{Until}(m\text{'s token})$, that
  is, \kw{AllOf} with one condition: a \kw{ForOne} waiter, FIFO by
  registration. \\
\kw{wait}$(\mathit{cv}, m)$ & Three Steps. (1) One \kw{Execute}, under $m$:
  append the thread ID $w$ to the Attribute $\mathit{waiters}/\mathit{cv}$
  and publish $m$'s token. (2) Wait on the per-waiter Channel
  $\mathit{cv}/w$. (3) \kw{lock}$(m)$ again. \\
\kw{signal}$(\mathit{cv})$ & One \kw{Execute}, under $m$: if
  $\mathit{waiters}/\mathit{cv}$ is not empty, remove its head $h$ and
  publish to $\mathit{cv}/h$. \\
\kw{broadcast}$(\mathit{cv})$ & As \kw{signal}, for every waiter in
  order. \\
\kw{cancel}$(t)$ & \kw{Execute} publishes to $\mathit{cancel}/t$. \\
\kw{recv}$(c)$ & \kw{WaitFor} returns
  $\kw{AnyOf}(c,\ \mathit{cancel}/\mathit{self})$. The environment's
  publications fill $c$. \\
\kw{sleep}$(d)$ & \kw{GoTo} the \lean{timed} Step (Enter). Its \kw{WaitFor}
  returns $\kw{AnyOf}(\mathit{Timer}(d),\ \mathit{cancel}/\mathit{self})$
  (Arm). \\
\kw{select} with \kw{after}$(d)$ & As \kw{sleep}, with the wait
  $\kw{AnyOf}(\mathit{Timer}(d),\ c,\ \mathit{cancel}/\mathit{self})$. Only
  the winner consumes. \\
\kw{emit}$(a)$ & An \kw{Execute} whose only external effect is $a$. \\
\bottomrule
\end{tabularx}
\caption{The translation, by construct (\lean{encodeIntent},
\lean{encodeWaitRequest}, \lean{timedStepDef} in
\file{Translation/Flow.lean}). Every Channel condition is \kw{ForOne}.}
\label{tab:mapping}
\end{table}

The translation uses \kw{GoTo}, \kw{GoToMany}, \kw{DeadEnd},
\kw{ForceComplete}, and \kw{ForceFail}, and nothing else: no
\kw{CancelSteps}, no \kw{GracefulComplete}, no RPCs, no SubFlows, no
Attribute locks. Every decision it returns is well formed for every input
(\lean{translate_execute_isWellFormed}).

\paragraph{Why the encodings are faithful.}
Each encoding rests on a property of the Dex core from
\S\ref{sec:dex-rules}.

\emph{Mutex.} A mutex is a Channel that holds one token exactly while the
mutex is unlocked. The Channel has only \kw{ForOne} waiters, so a wake
serves them in registration order, and a token that is present when a
waiter registers means that no earlier waiter was blocked. This gives
\Lang{}'s FIFO acquisition. An Attribute lock would not do: it is held for
one Step method, not across a wait.

\emph{Condition variable.} The waiter list is an Attribute. Reading and
writing it in one \kw{Execute} would normally be a lost-update hazard
(\S\ref{sec:dex-rules}); here it is safe because the \kw{Execute} runs
while its thread holds $m$, and no other thread touches the list. This is
where the mutex requirement enters: \Lang{} does not let a thread begin
such a statement without $m$, the translated Step has no such check, and
well-formedness says that the difference never shows. The wake goes to a
per-waiter Channel, so a signal reaches the dequeued waiter even if that
waiter has not registered its wait yet: the message stays queued. No
spurious wakeup is needed. The waiter list is an AttributeMap instance, so
the Step that reads it must select it in its load options; the model
applies the selection as the server does, and
\lean{readWaiters_snapshot_of_cvOperand} proves that the declared selection
is enough.

\emph{Cancellation.} A cancel is a message on the per-thread Channel
$\mathit{cancel}/t$. \kw{AnyOf} consumes only from the winning condition,
so a cancel never eats an input message or a join token. At registration no
timer is done, and the Channel conditions decide in declaration order, the
operation's own Channel before $\mathit{cancel}/\mathit{self}$. That is
\Lang{}'s rule that a cancel is observed only if the point would block. The
translation never cancels a Step. The only Step cancellation in a
translated Flow is the Flow close, after which nothing is observable except
effects already in flight, which \Lang{}'s asynchronous \kw{emit} allows.

\emph{Timers.} Every timed operation runs in the one Step type
\lean{timed}, with one timer at index 0. \Lang{}'s Enter is the commit that
creates the Step, so the timer number is the Step number. \Lang{}'s Arm is
the \kw{WaitFor} completion, so the deadline is the fire time. For equal
fire times Dex orders timers by encoded timer ID, which for these IDs is
numeric order on the Step number. A \rulename{TIMER-FIRE} prefix of size
$k$ is then $k$ consecutive \kw{fire} actions of \Lang{}.

\emph{Exit.} \kw{ForceComplete} drops a null output, so the translation
passes the tagged value $\{\texttt{exit}: v\}$, which is never null even
when $v$ is \Lang{}'s null. The projection reads $v$ back from the one
output of the closed Flow.
